随着生成式人工智能逐步进入图像、音频和视频生产场景，地方文旅账号开始借助 AI 生成画面、虚拟场景和超现实表达开展短视频传播。本文以抖音平台省级文旅账号发布的相关短视频为研究对象，经筛选与复核后形成 722 条深度分析样本，用以考察内容特征对用户参与行为的影响。研究采用内容分析、差异检验、多元回归、稳健性检验和预测性补充分析等方法，并以点赞、评论、收藏和转发等公开互动行为衡量用户参与。

研究发现，样本互动量呈现明显的长尾分布。单因素检验表明，不同内容变量在不同类别之间存在显著差异。多元回归结果显示，在控制其他因素后，AI 技术质感、视频风格、视频时长和发布距采集日天数是较为稳定的影响因素。相较于粗糙质感，普通质感和精良质感均能显著提升传播效果；抽象/超现实风格相较于纪实/写实风格具有显著正向影响；视频时长对互动表现具有显著负向影响，发布距采集日天数则显著促进互动积累。机器学习结果进一步表明，AI 技术质感、视频时长和视频风格具有较高的变量重要性。

本文认为，AI 技术本身并不会自动带来参与优势。对地方文旅账号而言，关键不在于是否使用 AI，或 AI 介入程度有多深，而在于 AI 是否最终转化为更高的技术质感、更鲜明的视觉风格和更符合平台传播逻辑的观看体验。本文结论可为地方文旅部门优化短视频内容生产和提升平台用户参与提供参考。
